\documentclass[10pt,a4paper]{article}
\usepackage[utf8]{inputenc}
\usepackage{amsmath}
\usepackage{amsfonts}
\usepackage{amssymb}
\usepackage{graphicx}
\title{Rust notes}
\author{Keith Reid}
\begin{document}
\maketitle{}
\section{Intro}
I have decided to use pdflatex and papers to write my ReadMe from now on.

\section{Build}
There are two ways to build:
\begin{itemize}
  \item naive fragile way is to write the main.rs:\\
    \texttt{\$ rustc main.rs}\\
    \texttt{\$ ./main}\\
    \texttt{\$ Hello, world!}\\

  \item but the slightly trickier more robust way is to use cargo build:\\
    \texttt{\$ cargo new hello\_cargo}\\
    \texttt{\$ cd hello\_cargo}\\
    ...which lets you do \texttt{\$ cargo build} or \texttt{\$ cargo run}
\end{itemize}

As the docs say, 

\begin{itemize}
  \item create a project using \texttt{\$ cargo new}
  \item build a project using \texttt{\$ cargo build}
  \item build and run a project in one step using \texttt{\$ cargo run}
  \item build a project without producing a binary using \texttt{\$ cargo check}
   \item Cargo stores binaries in the target/debug directory.
   \item at p50 Urwin takes us through the \#[test] attribute.
\end{itemize}

\end{document}